\worksheettitle
  {Цифры, числа и разряды}
  {\modulemark{M01} Учебный модуль \textbullet\ ориентир 15--20 минут}

\studentnameblock

\begin{checkpointbox}[Вход: проверьте, готовы ли вы начать]
  Выполните без подсказки.
  \begin{enumerate}[label=\textbf{\arabic*.}]
    \item Сколько цифр в записи числа \(505\)? \answerblank[16mm]
    \item Какая цифра стоит в разряде десятков числа \(684\)?
      \answerblank[16mm]
  \end{enumerate}
\end{checkpointbox}

\section{Одна цифра --- разные значения}

Рассмотрите записи и заполните таблицу.

\begin{center}
\renewcommand{\arraystretch}{1.45}
\begin{tabularx}{0.94\linewidth}{@{}>{\centering\arraybackslash}p{0.18\linewidth}>{\centering\arraybackslash}X>{\centering\arraybackslash}X@{}}
  \toprule
  \rowcolor{TableHeader}
  Число & Разряд цифры \(6\) & Значение цифры \(6\) \\
  \midrule
  \(6\) & единицы & \(6\) \\
  \(60\) & \answerblank[25mm] & \answerblank[20mm] \\
  \(600\) & \answerblank[25mm] & \answerblank[20mm] \\
  \(6\,000\) & \answerblank[25mm] & \answerblank[20mm] \\
  \bottomrule
\end{tabularx}
\end{center}

\begin{reflectionbox}[Главный вывод]
  Значение цифры зависит от её \answerblank[32mm], то есть от
  \answerblank[32mm].
\end{reflectionbox}

\section{Разбираем число по разрядам}

\levelbase

Рассмотрите число \(72\,649\). Разряды определяйте справа налево.

\begin{center}
\renewcommand{\arraystretch}{1.45}
\begin{tabular}{c|ccccc}
  \toprule
  цифра & 7 & 2 & 6 & 4 & 9 \\
  \midrule
  разряд & \answerblank[21mm] & \answerblank[21mm] & \answerblank[16mm] &
    \answerblank[16mm] & \answerblank[16mm] \\
  значение & \answerblank[18mm] & \answerblank[18mm] & \answerblank[14mm] &
    \answerblank[14mm] & \answerblank[14mm] \\
  \bottomrule
\end{tabular}
\end{center}

\begin{practicebox}[Самостоятельная проба]
  \begin{enumerate}[label=\textbf{\arabic*.}]
    \item В числе \(47\,382\) цифра \(3\) стоит в разряде
      \answerblank[28mm] и имеет значение \answerblank[20mm].
    \item В числе \(26\,415\) цифра \(6\) стоит в разряде
      \answerblank[28mm] и имеет значение \answerblank[20mm].
    \item В числе \(53\,257\) две цифры \(5\) имеют значения
      \answerblank[22mm] и \answerblank[22mm].
  \end{enumerate}
\end{practicebox}

\begin{checkpointbox}[Контрольная точка]
  Рассмотрите число \(44\,307\).
  \begin{enumerate}[label=\textbf{\arabic*.}]
    \item Значение первой слева цифры \(4\): \answerblank[22mm].
    \item Значение второй слева цифры \(4\): \answerblank[22mm].
    \item Почему одинаковые цифры имеют разные значения?
      \writinglines{1}
  \end{enumerate}
\end{checkpointbox}

\begin{reflectionbox}[Моя готовность]
  \choicebox\ Я называю и разряд, и значение цифры.

  \choicebox\ Я называю цифру, но пока путаю её значение.
\end{reflectionbox}

